\chapter{8-3编码器——5大变式详解}

\section{变式1：16-4编码器}

\begin{detailbox}[完整教学]
\textbf{【电路升级】}
8-3编码器：$2^3=8$个输入→3位输出。
16-4编码器：$2^4=16$个输入→4位输出。

\textbf{【设计思想】}
原理完全不变——每个输出位=所有编号中该位为1的输入的OR。

对于16-4编码器，输出$Y_3Y_2Y_1Y_0$，输入$I_0\sim I_{15}$。

\textbf{【逻辑推导】}
\begin{align}
Y_3 &= I_8+I_9+I_{10}+I_{11}+I_{12}+I_{13}+I_{14}+I_{15} \quad\text{(编号8-15的bit3=1)}\\
Y_2 &= I_4+I_5+I_6+I_7+I_{12}+I_{13}+I_{14}+I_{15} \quad\text{(bit2=1的编号)}\\
Y_1 &= I_2+I_3+I_6+I_7+I_{10}+I_{11}+I_{14}+I_{15}\\
Y_0 &= I_1+I_3+I_5+I_7+I_9+I_{11}+I_{13}+I_{15}
\end{align}

\textbf{【通用公式】}对于$2^n$-$n$编码器：
$$Y_k = \sum_{i=0}^{2^n-1} I_i \quad\text{当} i \text{的第}k\text{位为1时}$$

\textbf{【VHDL】}
\begin{lstlisting}
entity encoder_16to4 is
  port (I : in std_logic_vector(15 downto 0);
        Y : out std_logic_vector(3 downto 0));
end encoder_16to4;

architecture behavioral of encoder_16to4 is
begin
  process(I)
  begin
    case I is
      when "0000000000000001" => Y <= "0000";
      when "0000000000000010" => Y <= "0001";
      -- ... 中间省略 ...
      when "1000000000000000" => Y <= "1111";
      when others => Y <= "0000";
    end case;
  end process;
end behavioral;
\end{lstlisting}

\textbf{【更好的写法——用for循环查找】}
\begin{lstlisting}
process(I)
begin
  Y <= "0000";
  for i in 0 to 15 loop
    if I(i) = '1' then
      Y <= std_logic_vector(to_unsigned(i, 4));
    end if;
  end loop;
end process;
\end{lstlisting}

\textbf{【考试聚焦】}$2^n$-$n$编码器的通用公式。16-4编码器需要16个case分支。
\end{detailbox}


\section{变式2：优先级编码器（必考！）}

\begin{detailbox}[完整教学]

\textbf{【问题】}
普通编码器假设同一时刻只有一个输入为1。但当多个输入同时为1时（比如多个按键同时按下），普通编码器产生错误输出。需要\textbf{优先级}——编号大的输入优先。

\textbf{【设计思想】}
用if-else级联：从高优先级到低优先级依次判断。I7优先级最高，I0最低。

\textbf{【逻辑推导】}
\begin{center}
\begin{tabular}{|c|c|c|c|c|c|c|c|c||c|c|c|}
\hline
I7 & I6 & I5 & I4 & I3 & I2 & I1 & I0 & Y2 & Y1 & Y0 \\
\hline
1 & x & x & x & x & x & x & x & 1 & 1 & 1 \\
0 & 1 & x & x & x & x & x & x & 1 & 1 & 0 \\
0 & 0 & 1 & x & x & x & x & x & 1 & 0 & 1 \\
0 & 0 & 0 & 1 & x & x & x & x & 1 & 0 & 0 \\
0 & 0 & 0 & 0 & 1 & x & x & x & 0 & 1 & 1 \\
0 & 0 & 0 & 0 & 0 & 1 & x & x & 0 & 1 & 0 \\
0 & 0 & 0 & 0 & 0 & 0 & 1 & x & 0 & 0 & 1 \\
0 & 0 & 0 & 0 & 0 & 0 & 0 & 1 & 0 & 0 & 0 \\
0 & 0 & 0 & 0 & 0 & 0 & 0 & 0 & 0 & 0 & 0 \\
\hline
\end{tabular}
\end{center}
（x=任意值，不管它是0还是1）

\textbf{【VHDL——if-else级联】}
\begin{lstlisting}
entity priority_encoder is
  port (I : in std_logic_vector(7 downto 0);
        Y : out std_logic_vector(2 downto 0);
        GS : out std_logic);  -- Group Select: 有输入=1
end priority_encoder;

architecture behavioral of priority_encoder is
begin
  process(I)
  begin
    if I(7) = '1' then
      Y <= "111"; GS <= '1';
    elsif I(6) = '1' then
      Y <= "110"; GS <= '1';
    elsif I(5) = '1' then
      Y <= "101"; GS <= '1';
    elsif I(4) = '1' then
      Y <= "100"; GS <= '1';
    elsif I(3) = '1' then
      Y <= "011"; GS <= '1';
    elsif I(2) = '1' then
      Y <= "010"; GS <= '1';
    elsif I(1) = '1' then
      Y <= "001"; GS <= '1';
    elsif I(0) = '1' then
      Y <= "000"; GS <= '1';
    else
      Y <= "000"; GS <= '0';  -- 无有效输入
    end if;
  end process;
end behavioral;
\end{lstlisting}

\textbf{【为什么if-else实现优先级】}
if-else的语义：从上到下依次判断，第一个为真的分支被执行。
所以I7在第一个if → I7优先级最高。

\textbf{【综合成什么硬件】}
级联MUX链（优先级MUX树）。I7→I6→...→I0逐级选择。

\textbf{【GS信号的必要性】}
没有GS时，输入全0（无有效输入）和只有I0=1时，输出都是000——无法区分！
GS=1表示"有有效输入"。

\textbf{【Testbench——测试优先级】}
\begin{lstlisting}
-- 同时置I7和I0=1 → 应该输出111(I7优先)
I <= "10000001"; wait for 10 ns;
-- 输出应为111, GS=1 (I7优先级高于I0)
\end{lstlisting}

\begin{examfocus}
\textbf{必考！}优先级编码器是考试最高频的组合逻辑变式之一。
\begin{itemize}
  \item if-else级联 = 优先级递减的硬件
  \item GS信号区分"无输入"和"选择I0"
  \item 优先级从高到低排列if条件
\end{itemize}
\end{examfocus}

\end{detailbox}


\section{变式3：带使能端的编码器}

\begin{detailbox}[完整教学]

\textbf{【电路】}
在标准编码器基础上增加Enable(E)输入。E=1时编码器正常工作；E=0时所有输出为0（或高阻态）。

\textbf{【VHDL】}
\begin{lstlisting}
entity encoder_with_enable is
  port (I : in std_logic_vector(7 downto 0);
        E : in std_logic;
        Y : out std_logic_vector(2 downto 0));
end encoder_with_enable;

architecture behavioral of encoder_with_enable is
begin
  process(I, E)
  begin
    if E = '0' then
      Y <= "000";  -- 禁止输出
    else
      -- 正常编码逻辑
      case I is
        when "00000001" => Y <= "000";
        when "00000010" => Y <= "001";
        -- ...
        when others => Y <= "000";
      end case;
    end if;
  end process;
end behavioral;
\end{lstlisting}

\textbf{【设计思想】}
使能端是数字电路中常见的控制信号。在不使用时关断输出以节省功耗/避免总线冲突。

\end{detailbox}


\section{变式4：带GS（Group Select）输出的编码器}

\begin{detailbox}[完整教学]

\textbf{【问题】}
普通8-3编码器：当所有输入为0时，输出000。但当$I_0=1$时，输出也是000。如何区分这两种情况？

\textbf{【方案】}增加GS（Group Select/Valid）输出：
\begin{itemize}
  \item 任意输入=1 → GS=1（表示输出有效）
  \item 所有输入=0 → GS=0（表示输出无效/无输入选中）
\end{itemize}

\textbf{【VHDL】}
\begin{lstlisting}
GS <= '0' when I = "00000000" else '1';
-- 或
GS <= I(0) or I(1) or I(2) or I(3) or I(4) or I(5) or I(6) or I(7);
\end{lstlisting}

\textbf{【GS硬件】}一个8输入OR门——需要8个输入引脚。

\textbf{【考试常考】}GS信号的意义：解决编码器输出的二义性问题。

\end{detailbox}


\section{变式5：低电平有效的编码器}

\begin{detailbox}[完整教学]

\textbf{【电路变化】}
输入和输出都是\textbf{低电平有效}（0表示有效/选中，1表示无效）。
这在实际IC（如74LS148）中很常见——低电平有效更节省功耗，且与TTL电平兼容。

\textbf{【逻辑推导】}
低电平有效时：
\begin{itemize}
  \item 输入端：$\overline{I_0}, \overline{I_1}, ..., \overline{I_7}$（上划线表示低电平有效）
  \item 当$\overline{I_k}=0$且其他为1时 → 输出编码为k
  \item 用NAND门替代OR门来实现（因为低有效的或=高有效的与非）
\end{itemize}

\textbf{【VHDL】}
\begin{lstlisting}
entity encoder_active_low is
  port (
    I_n : in  std_logic_vector(7 downto 0);  -- 低有效
    Y_n : out std_logic_vector(2 downto 0);  -- 低有效
    GS_n: out std_logic  -- 低有效：0=有输入, 1=无输入
  );
end encoder_active_low;

architecture behavioral of encoder_active_low is
begin
  process(I_n)
  begin
    GS_n <= '1';  -- 默认：无有效输入
    Y_n <= "111";

    if I_n(7) = '0' then
      Y_n <= "000"; GS_n <= '0';  -- 7的编码=111, 反相=000
    elsif I_n(6) = '0' then
      Y_n <= "001"; GS_n <= '0';  -- 6=110, 反相=001
    elsif I_n(5) = '0' then
      Y_n <= "010"; GS_n <= '0';
    elsif I_n(4) = '0' then
      Y_n <= "011"; GS_n <= '0';
    elsif I_n(3) = '0' then
      Y_n <= "100"; GS_n <= '0';
    elsif I_n(2) = '0' then
      Y_n <= "101"; GS_n <= '0';
    elsif I_n(1) = '0' then
      Y_n <= "110"; GS_n <= '0';
    elsif I_n(0) = '0' then
      Y_n <= "111"; GS_n <= '0';  -- 0的编码=000, 反相=111
    end if;
  end process;
end behavioral;
\end{lstlisting}

\textbf{【关键理解】}
低有效编码：输入0有效→输出也是0=111(即7的二进制反码)。高有效和低有效编码器的输出值互为反码。

\textbf{【考试提示】}看清题目要求：高有效还是低有效？74LS148是低有效优先级编码器，这是常考器件。

\end{detailbox}


\section{编码器变式总结}

\begin{examfocus}[编码器5大变式速查]
\begin{center}
\begin{tabular}{|l|l|l|}
\hline
\textbf{变式} & \textbf{核心改动} & \textbf{VHDL关键} \\
\hline
16-4编码 & 规模扩大 & case 16分支/for循环 \\
优先级编码 & 允许同时多输入 & if-else级联(高→低) \\
带使能 & 增加E控制 & if E='0' then Y<="000" \\
带GS输出 & 增加有效标志 & GS<=not(I=全0) \\
低电平有效 & 0有效1无效 & 判断I(n)='0', 输出取反 \\
\hline
\end{tabular}
\end{center}
\end{examfocus}
